
\chapter{External Interfaces}
\label{sec:ei-2}
\label{chap:ei-2}

%%%%%%%%%%%%%%%%%%%%%%%%%%%%%%%%%%%%%%%%%%%%%%%%%%%%%%%%%%%%%%%%%%%%%%
\section{Introduction}
\label{sec:ei-intro}

This chapter
contains
calls used to
create \mpitermdef{generalized requests},
which allow users 
to create new nonblocking
operations with an interface similar to what is present in \mpi/.
These calls can be used to layer new functionality on top of \mpi/. 
Section~\ref{sec:ei-status} deals with setting the information found
in
\mpiarg{status}.  
This functionality is needed for generalized 
requests.

%%%%%%%%%%%%%%%%%%%%%%%%%%%%%%%%%%%%%%%%%%%%%%%%%%%%%%%%%%%%%%%%%%%%%%

\section{Generalized Requests}
\mpitermtitleindex{generalized requests}
\label{sec:ei-gr}

The goal of 
generalized requests 
is to allow users to define new
nonblocking operations.  Such an outstanding nonblocking operation is
represented by a (generalized) request.  A fundamental property of 
nonblocking operations is that 
\mpitermni{progress}\mpitermindex{progress!generalized requests}\mpitermindex{generalized requests!progress}\mpitermindex{semantics!generalized requests!progress}
toward the completion of this
operation occurs asynchronously, i.e., concurrently with normal program
execution.  Typically, this requires execution of code concurrently
with the execution of the user code, e.g., in a separate thread or in a
signal handler.  Operating systems provide a variety of mechanisms in
support of concurrent execution.
\MPI/ does not attempt to standardize or 
to replace these mechanisms: it is assumed
programmers who wish to define new asynchronous operations will use
the  mechanisms provided by the underlying operating system. 
Thus, the calls in this section only provide a means for
defining the effect of \MPI/ calls such as \mpifunc{MPI\_WAIT} or
\mpifunc{MPI\_CANCEL} when they apply to generalized requests, and for
signaling to \MPI/ the completion of a generalized operation. 

\begin{rationale}
It is tempting to also define an \MPI/ standard mechanism for
achieving concurrent execution of
user-defined nonblocking operations.  
However, it is difficult to define such a mechanism without
consideration of the specific mechanisms used in the operating system.
The Forum feels that concurrency mechanisms are a proper part
of the underlying operating system and should not be standardized by
\MPI/; the \MPI/ standard should only deal with the interaction of
such mechanisms with \MPI/.
\end{rationale}

For a regular request, the operation associated with the request is
performed by the \MPI/ implementation, and the operation
completes without intervention by the application.  For a generalized
request, the operation associated with the request is performed by the
application; therefore, the application must notify \MPI/ 
through a call to \mpifunc{MPI\_GREQUEST\_COMPLETE} when the
operation completes.
\MPI/ maintains the ``completion'' status of generalized requests.  Any
other request state has to be maintained by the user. 
 
A new generalized request is started with

\cdeclindex{MPI\_Request}%
\begin{mpi-binding}
    function_name("MPI_Grequest_start")
    parameter(
        "query_fn",
        "FUNCTION",
        func_type="MPI_Grequest_query_function",
        desc=r"callback function invoked when request status is queried",
    )
    parameter(
        "free_fn",
        "FUNCTION",
        func_type="MPI_Grequest_free_function",
        desc=r"callback function invoked when request is freed",
    )
    parameter(
        "cancel_fn",
        "FUNCTION",
        func_type="MPI_Grequest_cancel_function",
        desc=r"callback function invoked when request is cancelled",
    )
    parameter("extra_state", "EXTRA_STATE", desc=r"extra state")
    parameter("request", "REQUEST", direction="out", desc=r"generalized request")
\end{mpi-binding}

\begin{users}
Note that a generalized request is of the same type as regular requests, in C and Fortran.
\end{users}

The call starts a generalized request and returns a handle to it in
\mpiarg{request}.

The syntax and meaning of the callback functions are listed below.
All callback functions are passed the \mpiarg{extra\_state} argument
that was associated with the request by the starting call
\mpifunc{MPI\_GREQUEST\_START}; \mpiarg{extra\_state} can be used to maintain user-defined
state for the request.  

In 
C, the query procedure is

\cdeclindex{MPI\_Status}%
\begin{mpi-binding}
    callback()
    no_lis_binding()
    no_render("f08, f90")
    
    function_name("MPI_Grequest_query_function", "GREQUEST_QUERY_FUNCTION")
    
    parameter("extra_state", "EXTRA_STATE")
    parameter("status", "STATUS", direction="out")
    
    no_ierror()
    parameter("ierror", "ERROR_CODE", suppress="c_parameter")
\end{mpi-binding}

\par\noindent
in Fortran with the \code{mpi\_f08} module

\begin{mpi-binding}
    render("MPI_Grequest_query_function", "f08")
\end{mpi-binding}

\par\noindent
in Fortran with the \code{mpi} module and (deprecated) \code{mpif.h} include file

\begin{mpi-binding}
    render("MPI_Grequest_query_function", "f90")
\end{mpi-binding}

\par\noindent

The \mpiarg{query\_fn} function computes the status
that should be returned for the generalized request.
The status
also includes information about successful/unsuccessful cancellation of
the request (result to be returned by \mpifunc{MPI\_TEST\_CANCELLED}).

The \mpiarg{query\_fn} callback is invoked by the 
\mpifuncindex{MPI\_WAITANY}% 
\mpifuncindex{MPI\_WAITSOME}% 
\mpifuncindex{MPI\_WAITALL}% 
\mpifuncindex{MPI\_TESTANY}% 
\mpifuncindex{MPI\_TESTSOME}% 
\mpifuncindex{MPI\_TESTALL}%
\mpifuncindex{MPI\_WAIT}%
\mpifuncindex{MPI\_TEST}%
\mpiskipfunc{MPI\_\{WAIT$|$TEST\}\{$|$ANY$|$SOME$|$ALL\}} call that
completed the generalized request associated with this callback.  
The callback function
is also 
invoked by calls to \mpifunc{MPI\_REQUEST\_GET\_STATUS}, if the request is
complete when the call occurs.  In 
both
cases, the callback is passed a reference to the corresponding status variable
passed by the user to the \MPI/ call; the status set by the callback
function is returned by the \MPI/ call. 
If the user provided \mpiconst{MPI\_STATUS\_IGNORE} or
\mpiconst{MPI\_STATUSES\_IGNORE} to the \MPI/ procedure that causes
\mpiarg{query\_fn} to be called, then \MPI/ will pass
a valid status object to \mpiarg{query\_fn}, and this status will be
ignored upon return of the callback function.  
Note that \mpiarg{query\_fn}
is invoked only after 
\mpifunc{MPI\_GREQUEST\_COMPLETE} is called on the request; 
it may be invoked several times for
the same generalized request, e.g., if the user calls
\mpifunc{MPI\_REQUEST\_GET\_STATUS} several times for this request.
Note also that a call to 
\mpifuncindex{MPI\_WAITSOME}%
\mpifuncindex{MPI\_WAITALL}%
\mpifuncindex{MPI\_TESTSOME}%
\mpifuncindex{MPI\_TESTALL}%
\mpiskipfunc{MPI\_\{WAIT$|$TEST\}\{SOME$|$ALL\}} may cause multiple
invocations of
\mpiarg{query\_fn} callback functions, one for each
generalized request that is completed by the \MPI/ call.  The order of
these invocations is not specified by \mpi/.

In C, the free procedure is

\begin{mpi-binding}
    callback()
    no_lis_binding()
    no_render("f08, f90")
    
    function_name("MPI_Grequest_free_function", "GREQUEST_FREE_FUNCTION")
    
    parameter("extra_state", "EXTRA_STATE")
    
    no_ierror()
    parameter("ierror", "ERROR_CODE", suppress="c_parameter")
\end{mpi-binding}

\par\noindent
in Fortran with the \code{mpi\_f08} module

\begin{mpi-binding}
    render("MPI_Grequest_free_function", "f08")
\end{mpi-binding}

\par\noindent
in Fortran with the \code{mpi} module and (deprecated) \code{mpif.h} include file

\begin{mpi-binding}
    render("MPI_Grequest_free_function", "f90")
\end{mpi-binding}

\par\noindent
The \mpiarg{free\_fn} function is invoked to clean up user-allocated
resources when the generalized request is freed. 

The \mpiarg{free\_fn} callback is invoked by the 
\mpifuncindex{MPI\_WAITANY}%
\mpifuncindex{MPI\_WAITSOME}%
\mpifuncindex{MPI\_WAITALL}%
\mpifuncindex{MPI\_TESTANY}%
\mpifuncindex{MPI\_TESTSOME}%
\mpifuncindex{MPI\_TESTALL}%
\mpifuncindex{MPI\_WAIT}%
\mpifuncindex{MPI\_TEST}%
\mpiskipfunc{MPI\_\{WAIT$|$TEST\}\{$|$ANY$|$SOME$|$ALL\}} call that
completed the generalized request associated with this
callback. \mpiarg{free\_fn} is invoked after the call to
\mpiarg{query\_fn} for the same request.  However, if the \MPI/ call
completed multiple generalized requests, the order in which
\mpiarg{free\_fn} callback functions are invoked is not specified by
\MPI/.
 
The \mpiarg{free\_fn} callback 
is also invoked for generalized requests that are freed by a call to
\mpifunc{MPI\_REQUEST\_FREE} (no call to 
\mpifuncindex{MPI\_WAITANY}%
\mpifuncindex{MPI\_WAITSOME}%
\mpifuncindex{MPI\_WAITALL}%
\mpifuncindex{MPI\_TESTANY}%
\mpifuncindex{MPI\_TESTSOME}%
\mpifuncindex{MPI\_TESTALL}%
\mpifuncindex{MPI\_WAIT}%
\mpifuncindex{MPI\_TEST}%
\mpiskipfunc{MPI\_\{WAIT$|$TEST\}\{$|$ANY$|$SOME$|$ALL\}}
will occur for
such a request).  In this case, the callback
function will be called either in the \MPI/ call
\mpifunc{MPI\_REQUEST\_FREE}\mpicode{(request)}, or in the \MPI/ call
\mpifunc{MPI\_GREQUEST\_COMPLETE}\mpicode{(request)}, whichever happens 
last, i.e., 
in this case the actual freeing code is executed
as soon as both calls \mpifunc{MPI\_REQUEST\_FREE} and
\mpifunc{MPI\_GREQUEST\_COMPLETE} have occurred. 
The \mpiarg{request} is not deallocated until after
\mpiarg{free\_fn} completes.
Note that \mpiarg{free\_fn} will be 
invoked only once per request by a correct program. 

\begin{users}
Calling \mpifunc{MPI\_REQUEST\_FREE}\mpicode{(request)} will cause the
\mpiarg{request} handle to be set to \mpiconst{MPI\_REQUEST\_NULL}.
This handle to the generalized request is no longer valid.  However,
user copies of this handle are valid until after
\mpiarg{free\_fn} completes since \MPI/ does not deallocate the object
until then.  Since \mpiarg{free\_fn} is not
called until after \mpifunc{MPI\_GREQUEST\_COMPLETE}, the user copy of
the handle can be used to make this call.  Users should note that
\MPI/ will deallocate the object after \mpiarg{free\_fn}
executes.  At this point, user copies of the \mpiarg{request} handle no
longer point to a valid request.  \MPI/ will not set user copies to
\mpiconst{MPI\_REQUEST\_NULL} in this case, so it is up to the user to
avoid accessing this stale handle.  This is a special case in which \mpi/
defers deallocating the object until a later time that is known by
the user.
\end{users}

In C, the cancel procedure is

\begin{mpi-binding}
    callback()
    no_lis_binding()
    no_render("f08, f90")
    
    function_name("MPI_Grequest_cancel_function", "GREQUEST_CANCEL_FUNCTION")
    
    parameter("extra_state", "EXTRA_STATE")
    parameter("complete", "LOGICAL")
    
    no_ierror()
    parameter("ierror", "ERROR_CODE", suppress="c_parameter")
\end{mpi-binding}

\par\noindent
in Fortran with the \code{mpi\_f08} module

\begin{mpi-binding}
    render("mpi_grequest_cancel_function", "f08")
\end{mpi-binding}

\par\noindent
in Fortran with the \code{mpi} module and (deprecated) \code{mpif.h} include file

\begin{mpi-binding}
    render("mpi_grequest_cancel_function", "f90")
\end{mpi-binding}

The \mpiarg{cancel\_fn} function is invoked to start the cancelation of
a generalized request.
It is called by \mpifunc{MPI\_CANCEL}\mpicode{(request)}.
\MPI/ passes \mpiarg{complete} = \mpicode{true} to the callback function if 
\mpifunc{MPI\_GREQUEST\_COMPLETE} was already called on the request,
and \mpiarg{complete} = \mpicode{false} otherwise.

All callback functions return an error code.  
The code is passed back and dealt with as appropriate for the error
code by the \MPI/ procedure that invoked the callback function.  For
example, if error codes are returned then the error code returned by
the callback function will be returned by the \MPI/ procedure that
invoked the callback function.
In the case of
an 
\mpifuncindex{MPI\_WAITANY}%
\mpifuncindex{MPI\_TESTANY}%
\mpiskipfunc{MPI\_\{WAIT$|$TEST\}\{ANY\}} call that invokes both
\mpiarg{query\_fn} and \mpiarg{free\_fn}, the \MPI/ call will return
the error code returned by the last callback, namely
\mpiarg{free\_fn}.  If one or more of the requests in a call to
\mpifuncindex{MPI\_WAITSOME}%
\mpifuncindex{MPI\_WAITALL}%
\mpifuncindex{MPI\_TESTSOME}%
\mpifuncindex{MPI\_TESTALL}%
\mpiskipfunc{MPI\_\{WAIT$|$TEST\}\{SOME$|$ALL\}} failed,
then the \MPI/ call will return
\error{MPI\_ERR\_IN\_STATUS}. 
In such a case, if the \MPI/ call was
passed an array of statuses, then \MPI/ will return in each of the
statuses that correspond to a completed generalized request the error
code returned by the corresponding invocation of its \mpiarg{free\_fn}
callback function.  However, if the \MPI/ procedure was passed
\mpiconst{MPI\_STATUSES\_IGNORE}, then the individual error codes
returned by each callback functions will be lost.

\begin{users}
\mpiarg{query\_fn} must \emph{not} set the error field of
\mpiarg{status}
since \mpiarg{query\_fn} may be called by \mpifunc{MPI\_WAIT} or
\mpifunc{MPI\_TEST}, in which case the error field of \mpiarg{status}
should not change.  The \MPI/ library knows the ``context'' in which
\mpiarg{query\_fn} is invoked and can decide correctly when to put
the returned error code in the error field of \mpiarg{status}.
\end{users}

\cdeclindex{MPI\_Request}%
\begin{mpi-binding}
    function_name("MPI_Grequest_complete")
    parameter(
        "request",
        "REQUEST",
        direction="lis:inout,param:in",
        desc=r"generalized request",
    )
\end{mpi-binding}

The call informs \MPI/ that the operations represented by the generalized request
\mpiarg{request} are 
complete (see 
definitions in Section~\ref{terms:semantic}).
A call to \mpifunc{MPI\_WAIT}\mpicode{(request, status)} will return and a call to
\mpifunc{MPI\_TEST}\mpicode{(request, flag, status)} will return \mpiarg{flag} = \mpicode{true}
only after a call to \mpifunc{MPI\_GREQUEST\_COMPLETE} has declared that
these operations are complete.

\MPI/ imposes no restrictions on the code executed by the callback functions.
However, new nonblocking operations should be defined so that the general
semantic rules about \MPI/ calls such as \mpifunc{MPI\_TEST},
\mpifunc{MPI\_REQUEST\_FREE}, or \mpifunc{MPI\_CANCEL} still hold. 
Therefore, the
callback functions \mpiarg{query\_fn}, \mpiarg{free\_fn}, or
\mpiarg{cancel\_fn} should invoke blocking \MPI/ communication
calls only if the context is such that these calls are guaranteed to
return in finite time.  
Once \mpifunc{MPI\_CANCEL} is invoked, the cancelled operation
should complete in finite time, irrespective of the state of
other \MPI/ processes (the operation has acquired ``local'' semantics).  It
should either succeed, or fail without side-effects.  The user should
guarantee these properties for newly defined operations.  

\begin{implementors}
A call to \mpifunc{MPI\_GREQUEST\_COMPLETE} may unblock a blocked user
process/thread. The \MPI/ library should ensure that the blocked user
computation will resume.
\end{implementors}

\subsection{Examples}

\begin{example}
\Exindex{C}{User-defined reduce}{MPI\_Grequest\_start,MPI\_Grequest\_complete}%
This example shows the code for a user-defined reduce operation on an
\code{int} using
a binary tree: each nonroot \MPI/ process receives two messages, sums them, 
and sends them up.  We assume that no status is returned and that the
operation cannot be cancelled.

%%HEADER
%%LANG: C
%%SKIPELIPSIS
%%DECL: #include <pthread.h>
%%DECL: int query_fn(void *extra_state, MPI_Status *status);
%%DECL: int free_fn(void *extra_state);
%%DECL: int cancel_fn(void *extra_state, int complete);
%%DECL: void* reduce_thread(void *ptr);
%%ENDHEADER
\begin{lstlisting}[language={[MPI]C}]
typedef struct {
   MPI_Comm comm;
   int tag;
   int root;
   int valin;
   int *valout;
   MPI_Request request;
   } ARGS;


int myreduce(MPI_Comm comm, int tag, int root,
             int valin, int *valout, MPI_Request *request)
{
   ARGS *args;
   pthread_t thread;
   
   /* start request */
   MPI_Grequest_start(query_fn, free_fn, cancel_fn, NULL, request);
   
   args = (ARGS*)malloc(sizeof(ARGS));
   args->comm = comm;
   args->tag = tag;
   args->root = root;
   args->valin = valin;
   args->valout = valout;
   args->request = *request;
   
   /* spawn thread to handle request */
   /* The availability of the pthread_create call is system dependent */
   pthread_create(&thread, NULL, reduce_thread, args);
   
   return MPI_SUCCESS;
}

/* thread code */
void* reduce_thread(void *ptr) 
{
   int lchild, rchild, parent, lval, rval, val;
   MPI_Request req[2];
   ARGS *args;
   
   args = (ARGS*)ptr;
   
   /* compute left and right child and parent in tree; set 
      to MPI_PROC_NULL if does not exist  */
   /* code not shown */
   ...
     
   MPI_Irecv(&lval, 1, MPI_INT, lchild, args->tag, args->comm, &req[0]);
   MPI_Irecv(&rval, 1, MPI_INT, rchild, args->tag, args->comm, &req[1]);
   MPI_Waitall(2, req, MPI_STATUSES_IGNORE);
   val = lval + args->valin + rval;
   MPI_Send(&val, 1, MPI_INT, parent, args->tag, args->comm);
   if (parent == MPI_PROC_NULL) *(args->valout) = val;
   MPI_Grequest_complete((args->request));   
   free(ptr);
   return(NULL);
}

int query_fn(void *extra_state, MPI_Status *status)
{
   /* always send just one int */
   MPI_Status_set_elements(status, MPI_INT, 1);
   /* can never cancel so always true */
   MPI_Status_set_cancelled(status, 0);
   /* choose not to return a value for this */
   status->MPI_SOURCE = MPI_UNDEFINED;
   /* tag has no meaning for this generalized request */
   status->MPI_TAG = MPI_UNDEFINED;
   /* this generalized request never fails */
   return MPI_SUCCESS;
}


int free_fn(void *extra_state)
{
   /* this generalized request does not need to do any freeing */
   /* as a result it never fails here */
   return MPI_SUCCESS;
}


int cancel_fn(void *extra_state, int complete)
{
   /* This generalized request does not support cancelling.
      Abort if not already done.
      If done then treat as if cancel failed.*/
   if (!complete) {
     fprintf(stderr,
             "Cannot cancel generalized request - aborting program\n");
     MPI_Abort(MPI_COMM_WORLD, 99);
   }
   return MPI_SUCCESS;
}
\end{lstlisting}
\end{example}

%%%%%%%%%%%%%%%%%%%%%%%%%%%%%%%%%%%%%%%%%%%%%%%%%%%%%%%%%%%%%%%%%%%%%%
\section{Associating Information with Status}
\mpitermtitleindex{status!associating information}
\label{sec:ei-status}

\MPI/ supports several different types of requests besides those for
point-to-point operations,
this includes
\mpi/ calls for I/O and generalized requests.  It is desirable to allow
these calls to use the same request mechanism, which allows one to
wait or test on different types of requests.  However,
\mpifuncindex{MPI\_WAITANY}%
\mpifuncindex{MPI\_WAITSOME}%
\mpifuncindex{MPI\_WAITALL}%
\mpifuncindex{MPI\_TESTANY}%
\mpifuncindex{MPI\_TESTSOME}%
\mpifuncindex{MPI\_TESTALL}%
\mpifuncindex{MPI\_WAIT}%
\mpifuncindex{MPI\_TEST}%
\mpiskipfunc{MPI\_\{TEST$|$WAIT\}\{$|$ANY$|$SOME$|$ALL\}} returns a status
with information about the request.  With the generalization of
requests, one needs to define what information will be returned in the
status object.

Each \MPI/
call fills in the appropriate fields in the status object.  
Any unused field will
have an undefined value.  A call to
\mpifuncindex{MPI\_WAITANY}%
\mpifuncindex{MPI\_WAITSOME}%
\mpifuncindex{MPI\_WAITALL}%
\mpifuncindex{MPI\_TESTANY}%
\mpifuncindex{MPI\_TESTSOME}%
\mpifuncindex{MPI\_TESTALL}%
\mpifuncindex{MPI\_WAIT}%
\mpifuncindex{MPI\_TEST}%
\mpiskipfunc{MPI\_\{TEST$|$WAIT\}\{$|$ANY$|$SOME$|$ALL\}} can modify any of
the fields in the status object.  Specifically, it can modify fields
that are undefined.  The fields with meaningful values for a given
request are defined in the respective sections.

Generalized requests raise additional considerations.  Here, the user
provides the functions to deal with the request.  Unlike other \mpi/
calls, the user needs to provide the information to be returned in
the status.  The status argument is provided directly to the callback
function where the status needs to be set.  Users can directly set the
values in 3 of the 5 status values.  The count and cancel fields are
opaque.  To overcome this, 
these calls 
are provided:

\cdeclindex{MPI\_Status}%
\begin{mpi-binding}
    function_name("MPI_Status_set_elements")
    parameter(
        "status",
        "STATUS",
        direction="inout",
        desc=r"status with which to associate count",
    )
    parameter("datatype", "DATATYPE", desc=r"datatype associated with count")
    parameter(
        "count",
        "POLYXFER_NUM_ELEM",
        desc=r"number of elements to associate with status",
    )
\end{mpi-binding}

This procedure modifies the opaque part of \mpiarg{status} so calls
to \mpifunc{MPI\_GET\_ELEMENTS}
will return \mpiarg{count}.
Calls to
\mpifunc{MPI\_GET\_COUNT} will return a compatible value.

\begin{rationale}
The number of elements is set instead of the count because the former
can deal with 
a 
non-integer number of datatypes.
\end{rationale}


A subsequent call to \mpifunc{MPI\_GET\_COUNT} or \mpifunc{MPI\_GET\_ELEMENTS} must use a
\mpiarg{datatype} argument that has the same type signature as the
\mpiarg{datatype} argument that was used in the call to
\mpifunc{MPI\_STATUS\_SET\_ELEMENTS}.

\begin{rationale}
The requirement of matching type signatures for these calls is similar to the restriction that holds 
when 
\mpiarg{count} is set by a
receive operation: in that case, calls to
\mpifunc{MPI\_GET\_COUNT} and \mpifunc{MPI\_GET\_ELEMENTS} must use a
\mpiarg{datatype} with the same signature as the datatype used in the
receive call.
\end{rationale}

\cdeclindex{MPI\_Status}%
\begin{mpi-binding}
    function_name("MPI_Status_set_cancelled")
    parameter(
        "status",
        "STATUS",
        direction="inout",
        desc=r"status with which to associate cancel flag",
    )
    parameter(
        "flag",
        "LOGICAL",
        desc=r"if true, indicates request was cancelled",
        direction="in",
    )
\end{mpi-binding}

If \mpiarg{flag} is set to \mpicode{true} then a subsequent call to 
\mpifunc{MPI\_TEST\_CANCELLED} will also return \mpiarg{flag}\mpicode{ = true},
otherwise it will return 
\mpiarg{false}.  


While the \mpiconst{MPI\_SOURCE}, \mpiconst{MPI\_TAG}, and \mpiconst{MPI\_ERROR}
status values are directly accessible by the user,
for convenience in some contexts, users can also
modify them via the procedure calls described below.
Procedures for querying these fields from a
status object are defined in Section~\ref{subsec:pt2pt-status}.


\cdeclindex{MPI\_Status}%
\begin{mpi-binding}
    function_name("MPI_Status_set_source")
    parameter(
        "status",
        "STATUS",
        direction="inout",
        desc=r"status with which to associate source rank",
    )
    parameter(
        "source",
        "RANK",
        desc=r"rank to set in the \mpiconst{MPI\_SOURCE} field",
        direction="in",
    )
\end{mpi-binding}

Set the \mpiconst{MPI\_SOURCE} field in the \mpiarg{status} object to the provided \mpiarg{source} argument.

\cdeclindex{MPI\_Status}%
\begin{mpi-binding}
    function_name("MPI_Status_set_tag")
    parameter(
        "status",
        "STATUS",
        direction="inout",
        desc=r"status with which to associate tag",
    )
    parameter(
        "tag",
        "TAG",
        desc=r"tag to set in the \mpiconst{MPI\_TAG} field",
        direction="in",
    )
\end{mpi-binding}

Set the \mpiconst{MPI\_TAG} field in the \mpiarg{status} object to the provided \mpiarg{tag} argument.

\cdeclindex{MPI\_Status}%
\begin{mpi-binding}
    function_name("MPI_Status_set_error")
    parameter(
        "status",
        "STATUS",
        direction="inout",
        desc=r"status with which to associate error",
    )
    parameter(
        "err",
        "ERROR_CODE",
        desc=r"error to set in the \mpiconst{MPI\_ERROR} field",
        direction="in",
    )
\end{mpi-binding}

Set the \mpiconst{MPI\_ERROR} field in the \mpiarg{status} object to the provided \mpiarg{err} error code.

\begin{rationale}
These functions exist for convenience when using \mpi/ from languages
other than C and Fortran, where having a function in the \mpi/ library
with a known API reduces the need for utility code written in C.
\end{rationale}

\begin{users}
Users are advised not to reuse the status fields for values other than
those for which they were intended.  Doing so may lead to unexpected
results when using the status object.  For example, calling
\mpifunc{MPI\_GET\_ELEMENTS} may cause an error if the value is
out of range or it may be impossible to detect such an error.  The
\mpiarg{extra\_state} argument provided with a generalized request can
be used to return information that does not logically belong in
status.
Furthermore, modifying the values in a status set internally by \MPI/,
e.g., \mpifunc{MPI\_RECV}, may lead to unpredictable results and is
strongly discouraged.
\end{users}

